\chapter{Administrer une organisation}

\section{Tableau de bord}

Le tableau de bord de l'administrateur synthétise l'effectif, les clients liés,
la disponibilité suivie, les revenus réglés, les stations gérées et les
sessions de recharge. Les données sont limitées à l'organisation connectée et
à la période sélectionnée.

\manualscreenshot{admin-overview.png}
  {Vue d'ensemble d'une organisation}
  {fig:admin-overview}

Les actions rapides permettent d'ouvrir directement la création d'un employé,
la tarification, les stations ou les rapports. Avant de comparer deux valeurs,
vérifier la période affichée dans la bannière.

\section{Identité de l'organisation}

Les informations d'organisation comprennent notamment le nom, l'adresse de
contact, le téléphone et le logo. Le logo est réutilisé dans les vues destinées
aux employés et aux clients.

\begin{procedurebox}{Mettre à jour l'identité de l'organisation}
  \begin{enumerate}[leftmargin=*]
    \item Ouvrir les paramètres depuis le menu du profil.
    \item Accéder à la section consacrée à l'organisation.
    \item Modifier uniquement les champs autorisés.
    \item Importer, si nécessaire, un logo au format PNG, JPG ou WebP respectant
    la limite indiquée.
    \item Enregistrer puis vérifier l'aperçu dans le tableau de bord.
  \end{enumerate}
\end{procedurebox}

\section{Abonnement et facturation}

La page \emph{Subscription \& billing} sépare la facturation de l'organisation
des paiements effectués par les clients pendant une recharge. Elle présente le
statut commercial, la prochaine échéance, l'utilisation des quotas, le plan
actuel, les factures et les demandes de changement.

\manualscreenshotcompact{admin-billing.png}
  {Suivi commercial et capacité du plan de l'organisation}
  {fig:admin-billing}

\subsection{Cycle commercial}

\begin{longtable}{|L{3.1cm}|L{7.2cm}|L{4.8cm}|}
  \hline
  \rowcolor{ctmint}
  \textbf{Statut} & \textbf{Signification} & \textbf{Action attendue} \\ \hline
  \endfirsthead
  \hline
  \rowcolor{ctmint}
  \textbf{Statut} & \textbf{Signification} & \textbf{Action attendue} \\ \hline
  \endhead
  \texttt{trialing} &
  Organisation en période d'évaluation avec des quotas spécifiques. &
  Tester le service et choisir une offre avant l'échéance. \\ \hline
  \texttt{active} &
  Plan commercial en cours de validité. &
  Surveiller l'utilisation et la date de renouvellement. \\ \hline
  \texttt{past\_due} &
  Une échéance ou une facture nécessite une régularisation. &
  Vérifier les factures et contacter le gestionnaire commercial. \\ \hline
  \texttt{grace\_period} &
  Période transitoire après expiration de l'essai ou du plan. &
  Renouveler avant la date de suspension. \\ \hline
  \texttt{suspended} &
  Les fonctions opérationnelles sont suspendues ; l'accès commercial et au
  profil reste disponible. &
  Régulariser puis demander la restauration. \\ \hline
  \texttt{cancelled} &
  Souscription terminée. &
  Contacter la plateforme pour une nouvelle activation. \\ \hline
\end{longtable}

\subsection{Expiration et notifications}

La plateforme envoie des rappels aux administrateurs à 7, 3 et 1 jour avant
l'échéance. À l'expiration, l'organisation passe en période de grâce. Si aucune
régularisation n'est enregistrée avant la fin de ce délai, l'accès opérationnel
est suspendu automatiquement.

\expected{L'administrateur conserve l'accès nécessaire pour consulter la
situation commerciale et demander le renouvellement, sans pouvoir poursuivre
les opérations qui nécessitent un abonnement actif.}

\subsection{Valeurs par défaut de l'essai}

\begin{table}[H]
  \centering
  \caption{Quotas commerciaux par défaut de la période d'essai}
  \label{tab:trial-defaults}
  \begin{tabularx}{\textwidth}{|Y|C{3.2cm}|Y|}
    \hline
    \rowcolor{ctmint}
    \textbf{Paramètre} & \textbf{Valeur par défaut} & \textbf{Règle} \\ \hline
    Durée de l'essai & 14 jours & Configurable globalement de 7 à 90 jours. \\ \hline
    Bornes & 5 & Nombre maximal de stations intégrées pendant l'essai. \\ \hline
    Employés & 4 & Opérateurs et techniciens actifs ou invités ; administrateur
    exclu. \\ \hline
    Opérateurs & 2 & Sous-limite comprise dans le total des employés. \\ \hline
    Techniciens & 2 & Sous-limite comprise dans le total des employés. \\ \hline
    Délai de grâce & 7 jours & Configurable globalement de 1 à 30 jours. \\ \hline
  \end{tabularx}
\end{table}

\attention{Les valeurs réellement applicables sont celles affichées dans
l'interface. Le Super Admin peut modifier les valeurs par défaut pour les
nouvelles organisations. Les invitations en attente comptent dans les quotas.}
